hl|tc|Groups, Rings, and Fields
2.1) A Group is a set \(G\) equipped with a binary operation \( (\cdot) : G \times G \rightarrow G \) that associates an element \( a \cdot b \in G \ \forall \ a, b \in G \) such that,
li| \( \ a \cdot (b \cdot c) = (a \cdot b) \cdot c \ \ \forall a,b,c \in G \)         (G1 : Associativity)
li| \( \exists e \mid a \cdot e = a = e \cdot a \ \ \forall a \in G  \)                  (G2 : Existance of Identity)
li| \( \forall a \in G, \ \exists a^{-1} \mid a \cdot a^{-1} = e = a^{-1} \cdot a \)     (G3 : Existance of Inverse)
A group is called commutative (or abelian) if \( a \cdot b = b \cdot a \ \ \forall a,b \in G \).
A set \(M\) with an operation \( (\cdot) : M \times M \rightarrow M \) which satisfies (G1,G2) is called a Monoid.
Examples,
li| \( (Z, +) \) is a commutative group, but \( (Z, \times) \) is a monoid.
li| \( (Q, +) \) & \( (Q - \{0\}, \times) \) are commutative groups.
li| The set of bijections (permutations) \( f : S \rightarrow S \) & composition operation \( (\circ) \) is called the Symmetric Group of degree \(|S|\) (& order \(n!\)).
 